4) The identity w.r.t. $(\cdot)$ is known as identity element of the ring and is denoted as $1$.
5) If $(\cdot)$ is commutative then $R$ is called as commutative ring.
6) If the identity exists in $R$ then $R$ is called as ring with identity or ring with unity.

- Example: 1) $(\mathbb{Z}, +, \cdot)$, $(\mathbb{Q}, +, \cdot)$, $(\mathbb{R}, +, \cdot)$, $(\mathbb{C}, +, \cdot)$ all are commutative ring with unity.

2) $(\mathbb{Z}_n, +, \cdot)$ is a commutative ring with unity.

3) $\mathbb{Z}(\sqrt{p}) = \{a + b\sqrt{p} \mid a, b \in \mathbb{Z}, p \text{ is prime}\}$

Then $\mathbb{Z}(\sqrt{p})$ is a commutative ring with unity.

4) $\mathbb{Z}(i) = \{a + bi \mid a, b \in \mathbb{Z}\}$ is also a commutative ring with unity.

$\rightarrow$ Ring of Gaussian Integer.

5) $M_n(\mathbb{R}) \rightarrow$ set of all $n \times n$ matrices with real entries.
$(M_n(\mathbb{R}), +)$ is group.
$(M_n(\mathbb{R}), \cdot)$ is a semi-group.
$A \cdot (B + C) = AB + AC$
$(A + B) \cdot C = AC + BC$
$\therefore (M_n(\mathbb{R}), +, \cdot)$ is a non-commutative ring with unity.
